\section{Conclusion}
\label{sec:conclusion}

We set out to measure whether public Behaviour-Driven Development suites still function as oracles,
and whether that capability erodes. We built the instrument, proved that it can be driven end to end
in two ecosystems, pointed it at a pre-specified frame of public repositories --- and found that the
binding constraint is not statistical but infrastructural.

\subsection{Findings}

\begin{enumerate}
\item \textbf{The instrument works.} Mutation-based oracle measurement driven by the BDD runner itself
      is feasible in both Java (PIT with the Cucumber-JVM JUnit-Platform engine) and
      JavaScript/TypeScript (Stryker with the official Cucumber runner), with per-mutant coverage
      status and per-scenario attribution, which is what the differential-kill construct requires.
\item \textbf{The population is thin and skewed.} Of \Nframe{} repositories in the frame,
      \Nscreened{} were screened, \Neligible{} passed static eligibility, and \Ngreen{} produced a
      green, genuinely executing Gherkin baseline. The eligible set is dominated by BDD tooling ---
      reporting agents, step libraries, plugins and templates --- rather than by applications with a
      domain layer. Python is effectively absent from the eligible set, so the planned three-stratum
      design could not be realised.
\item \textbf{The boundary between metrics was confirmed as load-bearing.} Our own harness produced a
      direct local instance of cross-tool incomparability: the same domain logic yielded 19 mutants
      under PIT and 43 under Stryker. Any study that pools such numbers is measuring its tools, not
      its subject.
\item \textbf{Where an association could not be estimated, we say so.} The pre-specified confirmatory
      family is void under the protocol's own drop rule, because no repository was measurable
      end-to-end. We report that as \emph{not estimated}, which is a different statement from
      \emph{no association}, and we do not substitute a null result for an absent one.
\end{enumerate}

\subsection{Contributions revisited}

The artefact is a five-step protocol that an organisation can apply to its own suite, with mandatory
coverage conditioning, mandatory differential attribution, and explicit interpretation rules that
forbid the most common misreadings. Its practical value does not depend on the corpus-scale study
having succeeded: inside an organisation, the preconditions it needs --- a controllable build, a
domain layer, a green baseline --- are achievable, and the differential kill is then a routine
management metric rather than a research instrument.

The empirical contribution is a quantified account of why the obvious research design does not work on
the obvious population. We consider this worth publishing precisely because the naive design is
attractive and appears in the literature as an aspiration: a systematic mapping of 166 BDD papers
named the shortage of specification-quality metrics as a gap, and the natural response to that gap is
to measure quality on public repositories. Our attrition analysis shows what that response actually
encounters.

\subsection{Limitations we could not remove}

Three limitations are structural, and they are argued in full in Section~\ref{sec:threats}; we recall
them here only so the reader carries them past the results. No independent human raters, so the
judgement-coded indicators carry no reliability statistic. No build isolation, so the corpus is
environment-dependent and another host would likely yield a different $n$. And two snapshots, which
bound the longitudinal claim to a measured \emph{change} between snapshots rather than a rate of
erosion.

\subsection{Future work}

The study we could not run is the one worth running: a corpus of BDD suites drawn from organisations
that can control their build, services and data, measured longitudinally against business-rule change
events, with two independent raters coding the judgement-based indicators so that reliability can be
reported. Our instrument is designed to be reused there, our attrition analysis tells that study what
to budget, and the replication package contains the harnesses, the screening pipeline and the analysis
code needed to start from where this paper stopped rather than from the beginning.
